\chapter{REQUIREMENTS ANALYSIS}

\section{Dataset Requirements}
The model is trained on a merged, high-quality dataset comprising three primary sources, cleaned and formatted into a unified structure (Instruction, Input, Output) to support Supervised Fine-Tuning (SFT).

\begin{itemize}
    \item \textbf{Alpaca-18k:} A cleaned version of the general-purpose instruction-following dataset for maintaining conversational fluency.
    \item \textbf{Flytech-Coding:} Specialized technical data focused on programming logic and software engineering problems.
    \item \textbf{Refined Test/Train Sets:} Custom-curated \texttt{refined\_train.jsonl} and \texttt{refined\_test.jsonl} files ensure the model is evaluated on unseen, high-fidelity coding tasks.
\end{itemize}

\subsection{Data Processing and Masking}
\begin{itemize}
    \item \textbf{Merging \& Cleaning:} Deduplication and prompt template standardization were performed to ensure consistency.
    \item \textbf{Instruction Masking:} A custom tokenization script was implemented to set the \texttt{labels} of prompt tokens to $-100$. This ensures the loss is calculated only on the \texttt{code} completion, focusing the model's learning capacity on the output rather than the instructions.
    \item \textbf{Subsampling:} A random 5\% subset was extracted specifically for the \textbf{Optuna} hyperparameter tuning phase to reduce the computational cost of the 45-trial search.
\end{itemize}

\section{Hardware and Infrastructure Requirements}
Due to the computational demands of the Llama-2-7B architecture, a hybrid cloud infrastructure was utilized, leveraging both Kaggle and Google Colab environments.

\begin{table}[h]
    \centering
    \caption{Infrastructure Allocation and Hardware Resources}
    \begin{tabular}{|l|l|l|}
        \hline
        \textbf{Task} & \textbf{Platform} & \textbf{Hardware Resources} \\ \hline
        Data Preparation & Kaggle / Colab & CPU / T4 GPU (16GB VRAM) \\ \hline
        HPO (Optuna) & Kaggle CLI & Dual T4 GPUs / P100 \\ \hline
        Final Fine-tuning & Kaggle / Colab Pro & A100 (40GB) or L4 (24GB) \\ \hline
        Evaluation (HumanEval) & Kaggle & T4 GPU (Quantized Inference) \\ \hline
    \end{tabular}
\end{table}

\section{Software Requirements}
The implementation relies on a specialized ecosystem of Python libraries designed for memory-efficient training and standardized benchmarking.

\begin{itemize}
    \item \textbf{PEFT (Parameter-Efficient Fine-Tuning):} 
    Enables \textbf{LoRA (Low-Rank Adaptation)}. By freezing the 7-billion base parameters and training only small adapters (Rank $r=8$), the memory requirement for gradients is reduced by over 90\%. 
    
    

    \item \textbf{BitsAndBytes:} 
    Provides the backend for \textbf{4-bit Quantization (NF4)}. It includes "Double Quantization" and "Paged Optimizers," which allow the model to operate within 24GB VRAM by offloading optimizer states to CPU RAM during memory spikes.
    
    

    \item \textbf{Transformers \& TRL:} 
    Hugging Face's core libraries for model loading (\texttt{AutoModelForCausalLM}) and tokenization. The \texttt{Trainer} API was used to orchestrate the training loop and handle gradient accumulation.

    \item \textbf{Optuna:} 
    A hyperparameter optimization framework that utilizes the \textbf{Tree-structured Parzen Estimator (TPE)} to navigate the search space for learning rate, LoRA alpha, and dropout.

    \item \textbf{Weights \& Biases (W\&B):} 
    Used for real-time telemetry, monitoring validation loss curves, and logging hardware consumption across all 45 HPO trials.
\end{itemize}

\section{Quantization and LoRA Configuration}
To maintain efficiency while preserving 16-bit performance levels, QLoRA was implemented with the following optimal parameters derived from Trial 44:
\begin{itemize}
    \item \textbf{Precision:} 4-bit NormalFloat (nf4) with \texttt{torch.bfloat16} compute dtype.
    \item \textbf{LoRA Rank ($r$):} 8 (The low rank ensures the model doesn't overfit).
    \item \textbf{LoRA Alpha ($\alpha$):} 32 (A high scaling factor to amplify adapter influence).
    \item \textbf{Target Modules:} All linear layers (\texttt{q\_proj}, \texttt{k\_proj}, \texttt{v\_proj}, \texttt{o\_proj}, \texttt{gate\_proj}, \texttt{up\_proj}, \texttt{down\_proj}).
\end{itemize}

\section{Training and Evaluation Process}

\subsection{Hyperparameter Tuning Phase}
Automated HPO was conducted using Optuna. The process used a \textbf{Median Pruner} to terminate unpromising trials early. Trial 44 emerged as the best, achieving a validation loss of \textbf{0.5847} with a learning rate of $2.91 \times 10^{-5}$ and a cosine scheduler.

\subsection{Full Fine-tuning Phase}
The final training utilized the full merged dataset. To stabilize training, an \textbf{Effective Batch Size of 16} was used (Per-device batch 4 $\times$ Gradient Accumulation 4). **Gradient Checkpointing** was enabled to handle the 512-token sequence length on available hardware.

\subsection{Evaluation Pipeline (Pass@k)}
The model was evaluated using a three-step automated process:
\begin{enumerate}
    \item \textbf{Stochastic Inference:} Generating $n=5$ samples per problem using a temperature of $0.8$ to ensure diversity.
    \item \textbf{Regex Extraction:} A custom \texttt{re}-based script isolates code blocks from the model's natural language responses.
    \item \textbf{Unit Test Execution:} Utilizing the \texttt{human\_eval.execution} module to run generated code in a sandbox. The final success was measured using \textbf{Pass@1, Pass@2, and Pass@5} scores.
\end{enumerate}